% Options for packages loaded elsewhere
\PassOptionsToPackage{unicode}{hyperref}
\PassOptionsToPackage{hyphens}{url}
\documentclass[
]{article}
\usepackage{xcolor}
\usepackage{amsmath,amssymb}
\setcounter{secnumdepth}{-\maxdimen} % remove section numbering
\usepackage{iftex}
\ifPDFTeX
  \usepackage[T1]{fontenc}
  \usepackage[utf8]{inputenc}
  \usepackage{textcomp} % provide euro and other symbols
\else % if luatex or xetex
  \usepackage{unicode-math} % this also loads fontspec
  \defaultfontfeatures{Scale=MatchLowercase}
  \defaultfontfeatures[\rmfamily]{Ligatures=TeX,Scale=1}
\fi
\usepackage{lmodern}
\ifPDFTeX\else
  % xetex/luatex font selection
\fi
% Use upquote if available, for straight quotes in verbatim environments
\IfFileExists{upquote.sty}{\usepackage{upquote}}{}
\IfFileExists{microtype.sty}{% use microtype if available
  \usepackage[]{microtype}
  \UseMicrotypeSet[protrusion]{basicmath} % disable protrusion for tt fonts
}{}
\makeatletter
\@ifundefined{KOMAClassName}{% if non-KOMA class
  \IfFileExists{parskip.sty}{%
    \usepackage{parskip}
  }{% else
    \setlength{\parindent}{0pt}
    \setlength{\parskip}{6pt plus 2pt minus 1pt}}
}{% if KOMA class
  \KOMAoptions{parskip=half}}
\makeatother
\usepackage{graphicx}
\makeatletter
\newsavebox\pandoc@box
\newcommand*\pandocbounded[1]{% scales image to fit in text height/width
  \sbox\pandoc@box{#1}%
  \Gscale@div\@tempa{\textheight}{\dimexpr\ht\pandoc@box+\dp\pandoc@box\relax}%
  \Gscale@div\@tempb{\linewidth}{\wd\pandoc@box}%
  \ifdim\@tempb\p@<\@tempa\p@\let\@tempa\@tempb\fi% select the smaller of both
  \ifdim\@tempa\p@<\p@\scalebox{\@tempa}{\usebox\pandoc@box}%
  \else\usebox{\pandoc@box}%
  \fi%
}
% Set default figure placement to htbp
\def\fps@figure{htbp}
\makeatother
% definitions for citeproc citations
\NewDocumentCommand\citeproctext{}{}
\NewDocumentCommand\citeproc{mm}{%
  \begingroup\def\citeproctext{#2}\cite{#1}\endgroup}
\makeatletter
 % allow citations to break across lines
 \let\@cite@ofmt\@firstofone
 % avoid brackets around text for \cite:
 \def\@biblabel#1{}
 \def\@cite#1#2{{#1\if@tempswa , #2\fi}}
\makeatother
\newlength{\cslhangindent}
\setlength{\cslhangindent}{1.5em}
\newlength{\csllabelwidth}
\setlength{\csllabelwidth}{3em}
\newenvironment{CSLReferences}[2] % #1 hanging-indent, #2 entry-spacing
 {\begin{list}{}{%
  \setlength{\itemindent}{0pt}
  \setlength{\leftmargin}{0pt}
  \setlength{\parsep}{0pt}
  % turn on hanging indent if param 1 is 1
  \ifodd #1
   \setlength{\leftmargin}{\cslhangindent}
   \setlength{\itemindent}{-1\cslhangindent}
  \fi
  % set entry spacing
  \setlength{\itemsep}{#2\baselineskip}}}
 {\end{list}}
\usepackage{calc}
\newcommand{\CSLBlock}[1]{\hfill\break\parbox[t]{\linewidth}{\strut\ignorespaces#1\strut}}
\newcommand{\CSLLeftMargin}[1]{\parbox[t]{\csllabelwidth}{\strut#1\strut}}
\newcommand{\CSLRightInline}[1]{\parbox[t]{\linewidth - \csllabelwidth}{\strut#1\strut}}
\newcommand{\CSLIndent}[1]{\hspace{\cslhangindent}#1}
\setlength{\emergencystretch}{3em} % prevent overfull lines
\providecommand{\tightlist}{%
  \setlength{\itemsep}{0pt}\setlength{\parskip}{0pt}}
\usepackage{fancyhdr}
\pagestyle{fancy}
\fancyhf{}
\fancyfoot[C]{\footnotesize INTAKE --- NOT SUBMISSION READY --- AUTHOR METADATA DEFERRED}
\fancyfoot[R]{\thepage}
\setlength{\headheight}{14pt}
\usepackage{bookmark}
\IfFileExists{xurl.sty}{\usepackage{xurl}}{} % add URL line breaks if available
\urlstyle{same}
% fallback for those not using the hyperref driver hyperxmp:
\makeatletter
\@ifundefined{xmpquote}{\newcommand{\xmpquote}[1]{#1}}{}
\makeatother
\hypersetup{
  pdftitle={KEYSTONE-MPP-F1 -\/-\/- DRAFT PREAUTHORIZATION MANUSCRIPT},
  pdfauthor={Author metadata deferred},
  hidelinks,
  pdfcreator={LaTeX via pandoc}}

\ifXeTeX
\special{pdf:trailerid [ <2c22e8e75e21567a28aa50ae592a3e1e> <2c22e8e75e21567a28aa50ae592a3e1e> ]}
\fi
\ifPDFTeX
\pdftrailerid{}
\pdftrailer{/ID [ <2c22e8e75e21567a28aa50ae592a3e1e> <2c22e8e75e21567a28aa50ae592a3e1e> ]}
\fi
\ifLuaTeX
\pdfvariable trailerid {[ <2c22e8e75e21567a28aa50ae592a3e1e> <2c22e8e75e21567a28aa50ae592a3e1e> ]}
\fi

\title{KEYSTONE-MPP-F1 -\/-\/- DRAFT PREAUTHORIZATION MANUSCRIPT}
\author{Author metadata deferred}
\date{2026-08-30}

\begin{document}
\maketitle

\subsection{Title}\label{title}

KEYSTONE: A Draft Study of Dispute-Key Serviceability Distinct from
Ciphertext Availability in Encrypted AI Rollups

Status: \texttt{DRAFT\ /\ PRE-MANUSCRIPT\ /\ PRE-AUTHORIZATION}\\
System: \texttt{KEYSTONE-MPP-F1}\\
Date: \texttt{2026-08-29}\\
Canonical phase at draft time: \texttt{INTAKE}

\subsection{Abstract}\label{abstract}

Encrypted optimistic AI rollups can preserve ciphertext availability
while still failing to deliver timely dispute-key serviceability.
KEYSTONE studies that separation through a non-production
readiness-audit workflow assembled from established threshold-decryption
components rather than a new cryptographic primitive. The present draft
is bounded to three research-case claims: a property-separation
hypothesis (\texttt{C001}), an internal prototype and reproducibility
hypothesis (\texttt{C002}), and a broader solution-viability hypothesis
that still requires external evidence (\texttt{C003}). Existing tests,
simulations, figures, diagrams, and contract-boundary artifacts are
internal, analytic, preliminary, or exploratory evidence only. Novelty
remains unresolved, feasibility is unassessed in canonical state, and
confirmatory execution is not authorized. Accordingly, this shell
records admissible wording, planned sections, and explicit evidence
gaps; it does not make a publication-readiness claim.

\subsection{1. Introduction}\label{introduction}

The motivating problem is not ciphertext publication alone. A rollup
dispute may require authorized recovery of a record key or
threshold-decryption output within a deadline, and that requirement can
fail even when ciphertext bytes or receipts remain retrievable. The
current draft treats that gap as the core research question.

The manuscript remains within \texttt{FREEZE.md},
\texttt{docs/11\_PAPER\_BLUEPRINT.md}, and the canonical novelty case.
It does not claim a new threshold-decryption primitive, timing-free
accountability, adaptive selective-withholding resistance, unconditional
future availability, or production security.

The three manuscript-level claim anchors are:

\begin{itemize}
\tightlist
\item
  \texttt{C001}: study and formalize dispute-key availability as
  distinct from ciphertext availability under the frozen model, subject
  to unresolved strongest-prior-art closure;
\item
  \texttt{C002}: document the internal, non-production prototype and
  local reproducibility boundary;
\item
  \texttt{C003}: test the frozen analytic, simulated, prototype, and
  eventually external evidence package without wording above its current
  maturity ceiling.
\end{itemize}

\subsection{Methods}\label{methods}

The detailed methods architecture is organized across Sections 2--6.
Those sections remain a pre-authorization manuscript view of the
canonical novelty, protocol, analysis-plan, and implementation
artifacts.

The feasibility package now includes a non-executable confirmatory pilot
plan, machine-readable result contract, and a result-blind PC03
amendment. They reserve all nine \texttt{RID-*} result families,
separate the minimum short-paper core from the extended deadline lane,
prohibit reuse of historical outputs as confirmatory evidence, and leave
all observed outcomes missing. For three minimum synthetic cells---IID
baseline, stratified-versus-uniform at \texttt{s=8}, and selective
withholding at \texttt{w=11}---the draft amendment predeclares 32
primary and four reserve seed blocks, 4,096 draws per block,
result-blind precision rules, and deterministic seed schedules.
Correlated-domain analysis is
\texttt{EXCLUDED\_PENDING\_TRUTHFUL\_DOMAIN\_LABEL\_SOURCE}; its current
plots remain exploratory only. The distributed deadline family is
\texttt{EXCLUDED\_PENDING\_ENVIRONMENT\_PROFILE}. This design-readiness
package does not issue \texttt{GO}: it still needs independent methods
verification and separate execution authority.

\subsection{2. Background and Related
Work}\label{background-and-related-work}

\texttt{PENDING\ —\ NOVELTY\_UNRESOLVED.} This section will compare
KEYSTONE with data availability sampling, threshold encryption and
verifiable partial decryption, VSS/PVSS and DKG, proactive or dynamic
sharing, proof of retrievability, encrypted execution, threshold KMS,
and accountable liveness. The current novelty matrix rejects broad
primitive or first-system novelty. Until the independent, signed novelty
closure exists, only narrow property-separation and integration language
is admissible.

PeerDAS is the relevant contrast class for public or ciphertext-oriented
data availability rather than dispute-key serviceability (Ryan et al.
2024). Established secret-sharing, PVSS, proactive sharing, robust DPSS,
threshold-decryption, and DLEQ-style proof lineage are adopted building
blocks rather than claimable novelty (Shamir 1979; Schoenmakers 1999;
Herzberg et al. 1995; Yurek et al. 2023; Boneh et al. 2006; Chaum and
Pedersen 1993).

The adjacent encrypted-rollup, threshold-release, and
authorized-decryption design space is already populated by encrypted
optimistic-execution and threshold-release systems, including EigenAI,
ETHTID, Ferveo, vetKeys, BEAT-MEV, practical one-time-setup batched
threshold encryption, time-lock encrypted storage, context-dependent
threshold decryption, and consensus- authorized threshold decryption for
voting (Alves et al. 2026; Stengele et al. 2021; Bebel and Ojha 2022;
Cerulli et al. 2023; Bormet et al. 2025; Choudhuri et al. 2025; Agarwal
et al. 2025; Boneh et al. 2025; Liu et al. 2026). Accountability and
timing assumptions are also prior-art surfaces rather than
KEYSTONE-owned primitives (Boneh et al. 2023; Lewis-Pye et al. 2025).
Accordingly, the only currently admissible contribution language is the
narrower distinction between ciphertext availability and dispute-key
serviceability under the frozen model, plus the claim-bounded audit and
analysis surfaces that still require independent review.

\subsection{3. Model and Definitions}\label{model-and-definitions}

The pre-authorization formal source defines ciphertext availability,
valid-share registration, audit-time readiness, authorized
decryptability, deadline liveness, pre-authorization confidentiality,
and conditional dispute-key availability. Three constructive obligations
are currently preserved: (i) ciphertext availability does not imply
deadline-bounded dispute-key availability when only \texttt{t-1} record
contributions are serviceable; (ii) a routine audit may pass while a
target-aware selective withholder prevents an authorized dispute from
reaching \texttt{t} contributions; and (iii) a finite audit prefix
cannot guarantee unconditional future availability without temporal
assumptions. These are draft logical non-implication witnesses, not
prevalence, novelty, or production-security results. The paper does not
claim a complete pairwise property lattice, and all three obligations
remain subject to formal and novelty review.

\subsection{4. KEYSTONE Protocol}\label{keystone-protocol}

The authoritative pre-authorization protocol is
\texttt{research-case/03-design/protocol.md}. The manuscript-facing
protocol will cover epoch setup, record KEM/DEM handling, beacon-derived
canary challenges, context-bound partial decryptions and proofs, audit
verification, authorized dispute release, evidence records, recovery,
and domain-aware placement.

The current local package contains deterministic transcript and
signature fixtures that exercise audit-request and partial-response
interfaces. These are interoperability fixtures, not empirical evidence
or proof of production security. The current local audit API accepts
beacon and context inputs and derives the canary challenge internally,
removing the direct caller-supplied group-element path from routine
audit execution. This is a pre-authorization prototype boundary, not a
formal non-substitution proof or production-security claim, and it still
requires independent cryptographic review. The current draft treats
threshold decryption, DLEQ-style verification, and hash-to-curve usage
as adopted components rather than new cryptographic claims (Boneh et al.
2006, 2025; Chaum and Pedersen 1993; Faz-Hernandez et al. 2023).

\subsection{5. Analysis}\label{analysis}

The current analysis surface includes a static catastrophic-state
hypergeometric calculation and internal simulation models for IID,
correlated-domain, stratified-sampling, and selective-withholding
scenarios. Any static probability statement must be qualified as
one-audit reasoning under the declared fixed-ready-set and sampling
assumptions. Repeated-audit and deadline claims require explicit
temporal and synchrony assumptions.

Current boundaries are mandatory:

\begin{itemize}
\tightlist
\item
  temporal Markov outputs are exploratory only;
\item
  selective withholding is a negative finding and limitation;
\item
  deadline success remains blocked pending an accountable environment
  profile, synchronized trace design, and authorized distributed
  benchmark;
\item
  no analysis statement may exceed the claim ceilings encoded in the
  canonical research case.
\end{itemize}

\subsection{6. Implementation}\label{implementation}

The workspace contains a Python threshold-KEM/DLEQ prototype,
deterministic fixtures, local experiment scripts, and a Solidity
bulletin-board boundary. Local cryptographic timings and contract
test-body gas measurements are implementation artifacts only. They do
not establish production throughput, deployment cost, audited
cryptographic security, or independent reproducibility.

Two fail-closed pre-authorization engineering receipts currently bind
local source/configuration hashes to observed test outcomes. The crypto
receipt maps 22 pytest nodes and two vector-freshness checks; the
contract receipt maps 27 Foundry cases, including 14 deterministic
boundary tests, three fuzz properties, eight test-body gas operations,
and two stateful invariants. All mapped cases passed in the current
same-workspace runs. The receipts explicitly remain non-empirical,
non-independent engineering QA; they are not canonical scientific
results, cryptographic or smart-contract audits, or production
validation. An editable pre-authorization QA coverage table is generated
from the two receipts and their hash sidecars at
\texttt{paper/tables/preauthorization\_engineering\_qa.csv} for paper
drafting, but it is not a canonical T-series result table and cannot be
used to promote any scientific claim or raise the maturity or supported
scope of \texttt{C002}.

Production DKG/PVSS, audited cryptographic libraries, distributed
deployment, and independently controlled reproduction remain outside the
current internal prototype boundary.

\subsubsection{6.1 Draft protocol and method display
map}\label{draft-protocol-and-method-display-map}

The following diagram callouts bind the current editable sources and
rendered derivatives to their intended manuscript roles. They
communicate definitions, protocol structure, methods, and limitations;
they are not result evidence and do not raise the maturity of
\texttt{C001}, \texttt{C002}, or \texttt{C003}. Final placement remains
subject to the target venue, final-size review, and accountable-human
approval.

\begin{figure}
\centering
\pandocbounded{\includegraphics[keepaspectratio,alt={D1. System architecture separating ciphertext publication, policy coordination, readiness auditing, authorized release, threshold combination, private re-execution, and the dispute verdict.}]{diagrams/01_system_architecture.png}}
\caption{D1. System architecture separating ciphertext publication,
policy coordination, readiness auditing, authorized release, threshold
combination, private re-execution, and the dispute verdict.}
\end{figure}

\textbf{D1 --- System architecture.} Ciphertext availability, routine
canary auditing, and authorized dispute-key release are shown as
separate paths. The diagram is a protocol map for
\texttt{C002}/\texttt{C003}, not evidence that each path is externally
available or production-secure.

\begin{figure}
\centering
\pandocbounded{\includegraphics[keepaspectratio,alt={D2. Three constructive property-separation witnesses, each routing a held condition and a failed condition to one bounded non-implication.}]{diagrams/02_property_separation.png}}
\caption{D2. Three constructive property-separation witnesses, each
routing a held condition and a failed condition to one bounded
non-implication.}
\end{figure}

\textbf{D2 --- Constructive property separation.} Exactly three
witnesses support the draft analytic obligations in \texttt{C001}:
ciphertext availability does not imply deadline-bounded DKA; routine
audit acceptance does not unconditionally imply targeted dispute
success; and a finite audit prefix does not guarantee future DKA. This
is not a complete pairwise property lattice.

\begin{figure}
\centering
\pandocbounded{\includegraphics[keepaspectratio,alt={D3. Eight-step readiness-audit sequence from committed epoch data and beacon sampling through canary proof verification and evidence emission.}]{diagrams/03_audit_sequence.png}}
\caption{D3. Eight-step readiness-audit sequence from committed epoch
data and beacon sampling through canary proof verification and evidence
emission.}
\end{figure}

\textbf{D3 --- Readiness-audit sequence.} The routine path uses fresh,
context-bound canaries and classifies valid, invalid, and missing
responses at the declared deadline. It does not exercise production
ciphertext.

\begin{figure}
\centering
\pandocbounded{\includegraphics[keepaspectratio,alt={D4. Eight-step authorized-dispute sequence from challenge validation through confidential partial release, threshold combination, re-execution, and verdict.}]{diagrams/04_dispute_sequence.png}}
\caption{D4. Eight-step authorized-dispute sequence from challenge
validation through confidential partial release, threshold combination,
re-execution, and verdict.}
\end{figure}

\textbf{D4 --- Authorized dispute sequence.} The diagram records the
intended policy and threshold-release ordering; it does not establish
deployed confidential transport, deadline success, or production
authorization.

\begin{figure}
\centering
\pandocbounded{\includegraphics[keepaspectratio,alt={D5. Epoch, audit, and dispute state machines with explicit successful, failed, expired, and retired terminal paths.}]{diagrams/05_state_machines.png}}
\caption{D5. Epoch, audit, and dispute state machines with explicit
successful, failed, expired, and retired terminal paths.}
\end{figure}

\textbf{D5 --- State machines.} Allowed transitions make terminal and
failure paths explicit for implementation review. The current contract
is still a non-production boundary with admin-adjudicated off-chain
proof outcomes.

\begin{figure}
\centering
\pandocbounded{\includegraphics[keepaspectratio,alt={D6. Six adversary or fault classes mapped through KEYSTONE controls to six residual assumptions or risks.}]{diagrams/06_threat_model.png}}
\caption{D6. Six adversary or fault classes mapped through KEYSTONE
controls to six residual assumptions or risks.}
\end{figure}

\textbf{D6 --- Threat model and residual risk.} Each control is paired
with a residual assumption. In particular, domain labels may be
dishonest, audits do not prove future cooperation, and the
pre-authorization collusion threshold remains an assumption rather than
an observed result.

\begin{figure}
\centering
\pandocbounded{\includegraphics[keepaspectratio,alt={D7. Four fault domains with one offline domain and contrasting uniform versus domain-stratified sample routes.}]{diagrams/07_sampling_domains.png}}
\caption{D7. Four fault domains with one offline domain and contrasting
uniform versus domain-stratified sample routes.}
\end{figure}

\textbf{D7 --- Sampling and failure domains.} This is a method
illustration under truthful domain labels. It does not claim universal
superiority of stratification or field prevalence of the illustrated
outage.

\begin{figure}
\centering
\pandocbounded{\includegraphics[keepaspectratio,alt={D8. Reproducibility pipeline from frozen configuration through tests, simulation, data products, figures, claim mapping, manuscript assembly, and publication gates.}]{diagrams/08_experiment_pipeline.png}}
\caption{D8. Reproducibility pipeline from frozen configuration through
tests, simulation, data products, figures, claim mapping, manuscript
assembly, and publication gates.}
\end{figure}

\textbf{D8 --- Experiment-to-manuscript pipeline.} The flow separates
configuration, execution, versioned outputs, displays, claim binding,
and review gates. A clean local pipeline does not substitute for
authorized confirmatory analysis, independent reproduction, or human
submission approval.

\subsection{Results}\label{results}

The present results surface is an explicitly preliminary inventory. It
does not replace missing authorized \texttt{05-analysis} evidence or the
still-missing canonical primary-results ledger.

The temporary Markdown/LaTeX preliminary-table bundle has a fail-closed
same-workspace lineage receipt that binds its five input files,
exporter, selected fields and row identities, section mapping, and both
outputs. This improves draft reproducibility only; it is not a canonical
T-series result, scientific verification, or independent reproduction.

A second-agent, same-host isolated-copy mechanical rerun subsequently
executed the frozen ten-command package check under tool-level offline
flags. The first fresh attempt failed closed when Foundry rejected a
non-boolean offline setting; its hash-bound failure receipt and logs
were preserved. After the environment contract was corrected and
independently re-reviewed, a new operator completed all ten commands
with zero return codes, unchanged source and isolated-copy inventories,
and a valid \texttt{MECHANICAL\_PASS} receipt. This establishes bounded
same-host mechanical replay only: it used the same machine and
potentially the same tool caches, did not enforce network isolation at
the kernel level, did not execute the distributed deadline study, and is
neither external reproduction nor independent scientific verification,
institutional approval, venue approval, or evidence that changes the
current claim-maturity ceilings.

\subsection{7. Evaluation}\label{evaluation}

Only a limited set of result classes is admissible for drafting:

\begin{itemize}
\tightlist
\item
  analytic static catastrophic false-accept and detection calculations;
\item
  deterministic correctness and interoperability fixtures;
\item
  internal preliminary IID, correlated-domain, stratified-sampling, and
  selective-withholding simulation outputs;
\item
  internal local cryptographic timings and contract test-body gas
  measurements.
\end{itemize}

The current figures F1--F5 and \texttt{paper/tables/t1\_t8\_package.*}
are draft preauthorization display assets. They are not canonical
confirmatory tables or publication evidence. Missing or blocked
evaluation outputs include the canonical \texttt{05-analysis} primary
ledger, authorized confirmatory robustness and negative-result analyses,
deadline Figures F6--F7, the planned recovery-timeline Figure F8,
independent reproduction, and external validation for \texttt{C003}. The
canonical \texttt{05-analysis/results/exploratory-findings.csv} now
records the four temporal Markov rows under distinct \texttt{EXPLORE-*}
identifiers. It preserves the simulated source maturity but caps claim
support at \texttt{V0\ ASSERTED}: each row is explicitly post-hoc,
preauthorization-only, unauthorized, non-independent, and descriptive
rather than confirmatory. Its conditional denominator is the number of
trials ending in the catastrophic state, not all simulated trials. Draft
\texttt{04-data} provenance and evidence-status ledgers hash-bind the
source inputs; the deadline result remains explicitly \texttt{BLOCKED}
with no executed provenance row. These ledgers improve traceability only
and do not activate scientific support edges or advance the canonical
\texttt{INTAKE} phase.

The canonical draft \texttt{05-analysis/results/negative-findings.csv}
now records only the interval-separated selective-withholding gap
already present in the frozen simulation (\texttt{NEG-SW-11} through
\texttt{NEG-SW-14}). Each row remains a preauthorization simulated
negative finding, unauthorized, non-independent, and capped at
\texttt{V0\ ASSERTED}. This closes a traceability gap but is not an
authorized confirmatory result and does not provide deadline,
production, external, field, adaptive-adversary, or universal-security
evidence.

The canonical draft
\texttt{05-analysis/results/robustness-and-boundaries.csv} now preserves
all 17 frozen IID outage cells and all 12 frozen matched-seed sampling
policy cells as a deterministic traceability ledger. Every row is
explicitly \texttt{DRAFT\_ROBUSTNESS\_BOUNDARY\_ONLY}, preauthorization,
simulated, unauthorized, non-independent, and capped at
\texttt{V0\ ASSERTED}. Correlated-domain output is excluded because its
current result-ID binding is entangled with the quarantined Markov
asset; Markov, selective-withholding, and deadline rows are also
excluded from this ledger. The file therefore does not establish
confirmatory robustness, policy optimality, production behavior,
deadline performance, generality, or external validity.

The future confirmatory outputs are predeclared in
\texttt{02-feasibility/pilot-run-contract.csv}. The minimum core
comprises the separation witness, exact static result, deterministic
crypto and contract boundaries, and the limitation-bearing
selective-withholding family. IID and stratification are supporting
qualification families; correlated-domain analysis is excluded pending a
truthful domain-label source. The included draft counts and stream
schedule are recorded in \texttt{03-design/pc03-prospective-counts.csv}
and \texttt{03-design/pc03-seed-schedule.csv}; these are design inputs,
not results. The distributed deadline result remains an extended
full-paper family and must be excluded, rather than imputed, if its
environment and precision gates remain open.

\subsubsection{7.1 Preliminary quantitative display
map}\label{preliminary-quantitative-display-map}

The five available quantitative figures are included below only to
freeze their drafting roles and captions. Their source data, generator,
SVG/PNG hashes, and automated visual-QA status are recorded in the
source and figure manifests. They remain preliminary and must be
replaced or re-authorized from canonical analysis ledgers before
submission.

\begin{figure}
\centering
\pandocbounded{\includegraphics[keepaspectratio,alt={F1. Static one-audit detection bound under uniform sampling as the sample size changes.}]{prototype/results/figures/figure_1_theoretical_detection_bound.png}}
\caption{F1. Static one-audit detection bound under uniform sampling as
the sample size changes.}
\end{figure}

\textbf{F1 --- Static detection bound.} Analytic one-audit probability
under the declared fixed catastrophic state and uniform-sampling
assumptions; not a temporal availability or field-frequency result.

\begin{figure}
\centering
\pandocbounded{\includegraphics[keepaspectratio,alt={F2. Internal IID outage sweep showing reconstruction and audit-pass estimates with Wilson intervals.}]{prototype/results/figures/figure_2_iid_failure_sweep.png}}
\caption{F2. Internal IID outage sweep showing reconstruction and
audit-pass estimates with Wilson intervals.}
\end{figure}

\textbf{F2 --- IID failure sweep.} Internal preliminary simulation with
the configured scenario grid and uncertainty calculation; not
confirmatory or externally validated evidence.

\begin{figure}
\centering
\pandocbounded{\includegraphics[keepaspectratio,alt={F3. Internal correlated-domain outage sweep under declared truthful domain labels.}]{prototype/results/figures/figure_3_domain_diversity.png}}
\caption{F3. Internal correlated-domain outage sweep under declared
truthful domain labels.}
\end{figure}

\textbf{F3 --- Correlated-domain behavior.} Internal model-bounded
simulation; domain identity correctness and real outage structure remain
unverified assumptions.

\begin{figure}
\centering
\pandocbounded{\includegraphics[keepaspectratio,alt={F4. Selective-withholding negative-result figure showing that routine audit acceptance can coexist with failed targeted dispute reconstruction.}]{prototype/results/figures/figure_4_selective_withholding_gap.png}}
\caption{F4. Selective-withholding negative-result figure showing that
routine audit acceptance can coexist with failed targeted dispute
reconstruction.}
\end{figure}

\textbf{F4 --- Selective-withholding limitation.} Preserved negative
result that blocks an unconditional audit-to-dispute-success inference.

\begin{figure}
\centering
\pandocbounded{\includegraphics[keepaspectratio,alt={F5. Matched internal comparison of uniform and domain-stratified sampling under the configured conditioned catastrophic states.}]{prototype/results/figures/figure_5_sampling_strategy.png}}
\caption{F5. Matched internal comparison of uniform and
domain-stratified sampling under the configured conditioned catastrophic
states.}
\end{figure}

\textbf{F5 --- Sampling-strategy comparison.} A matched, model-specific
internal comparison. It does not establish universal dominance, external
validity, or production benefit.

\subsubsection{7.2 Draft evidence-bound table
map}\label{draft-evidence-bound-table-map}

The editable T1--T8 package is now available in
\texttt{paper/tables/t1\_t8\_package.md} and
\texttt{paper/tables/t1\_t8\_package.tex}. Every table remains draft /
pre-authorization only and inherits the canonical ceiling of the
underlying claim or evidence source.

\begin{itemize}
\tightlist
\item
  \texttt{T1} bounds the strongest-prior-art matrix and does not clear
  novelty.
\item
  \texttt{T2} is a design-comparator registry only and makes no
  measured-superiority claim.
\item
  \texttt{T3} records frozen local conditions and the still-missing
  \texttt{RID-C003-DEADLINE-001} distributed profile.
\item
  \texttt{T4} displays preauthorization simulated results parsed from
  canonical JSON outputs only, capped at \texttt{V0\ ASSERTED} for
  manuscript claims.
\item
  \texttt{T5} is a planned mechanism-isolation registry; every row
  remains \texttt{MISSING\_NOT\_EXECUTED}.
\item
  \texttt{T6} presents the canonical robustness ledger plus separately
  classified exact and exploratory temporal checks;
  selective-withholding negative findings remain in \texttt{T8}.
\item
  \texttt{T7} limits feasibility evidence to local timings and Foundry
  gas while leaving distributed, operator, network, and external rows
  blocked.
\item
  \texttt{T8} preserves selective withholding, synchrony limits,
  truthful-domain-label dependency, and missing external validation as
  explicit negative or unresolved findings.
\end{itemize}

\subsection{8. Discussion}\label{discussion}

The current safe interpretation is that ciphertext availability and
dispute-key serviceability should not be treated as interchangeable;
routine readiness audits may characterize present serviceability only
under stated assumptions; failure-domain modeling and sampling semantics
affect the meaning of reported probabilities; and local prototype
success does not establish external reproducibility, field readiness, or
production security.

Stronger interpretation must await reconciled prior art, authorized
analysis, and claim-matched independent evidence.

\subsection{9. Limitations and Future
Work}\label{limitations-and-future-work}

\begin{itemize}
\tightlist
\item
  Novelty remains \texttt{UNRESOLVED}; the broad historical claim is not
  submission-safe.
\item
  The canonical feasibility decision remains \texttt{UNASSESSED}; the
  draft feasibility package recommends a blocked posture and only a
  conditional \texttt{PILOT\_FIRST} path after prerequisites clear.
\item
  Current stochastic outputs are internal and mostly pre-authorization.
\item
  Temporal dependence is exploratory; selective withholding is a
  preserved negative result.
\item
  The distributed deadline benchmark and independent external validation
  are absent.
\item
  Independent threshold-cryptography and distributed-systems reviews are
  absent.
\item
  No venue portfolio, current reporting checklist, hermetic PDF package,
  rendered-page review, or accountable human approval exists.
\end{itemize}

\subsection{10. Conclusion}\label{conclusion}

\texttt{PENDING\ —\ EVIDENCE-BOUND\ SYNTHESIS.} The conclusion may
restate the distinction between ciphertext availability and dispute-key
serviceability and the conditional nature of any audit assurance. It
must not imply future-proof availability, external validity, production
readiness, or publication readiness unless the corresponding canonical
gates later pass.

\subsection{Data and Code
Availability}\label{data-and-code-availability}

Current local source assets include \texttt{prototype/src/keystone/*},
\texttt{prototype/tests/*}, \texttt{prototype/results/*},
\texttt{contracts/src/*}, \texttt{contracts/test/*},
\texttt{contracts/gas\_report.csv}, \texttt{paper/test\_vectors.json},
\texttt{paper/signature\_test\_vectors.json}, \texttt{diagrams/*}, and
the two \texttt{*/results/engineering\_qa/*-preauth-receipt.json}
bundles. Availability does not imply evidence authorization or final
reproducibility. The manuscript assembly inventory is tracked in
\texttt{07-manuscript/source-manifest.json}. Draft source-to-result
lineage is recorded in \texttt{04-data/provenance-manifest.csv} and
\texttt{04-data/evidence-status.csv}; those files retain the
preauthorization and same-workspace limitations of their inputs. The
derived temporal ledger at
\texttt{05-analysis/results/exploratory-findings.csv} is reproducible
from the hash-bound Markov dataset through
\texttt{scripts/export\_exploratory\_findings.py}, but is quarantined
from confirmatory, novelty, feasibility, and external-validation claims.
The derived robustness ledger at
\texttt{05-analysis/results/robustness-and-boundaries.csv} is
reproducible from the hash-bound IID and sampling-policy datasets
through \texttt{scripts/export\_robustness\_boundaries.py}; it remains a
draft traceability surface and cannot substitute for the missing primary
ledger, authorized confirmatory analysis, independent reproduction, or
external validation. The non-executable future-run contract is recorded
in \texttt{02-feasibility/pilot-plan.md} and
\texttt{02-feasibility/pilot-run-contract.csv}; it contains no collected
outcome and grants no authorization. The PC03 amendment, count contract,
and seed schedule under \texttt{03-design/} are likewise result-blind,
non-executable planning artifacts pending independent methods
verification.

\subsection{Declarations and Pending Human
Inputs}\label{declarations-and-pending-human-inputs}

Draft author metadata inputs have been supplied for prospective authors,
and the accountable human reported no institution, employer, funding
agreement, NDA, venue policy, or export-control conflict with the
current charter or study profile. At the author's request, final author
order, corresponding-author designation, accountable metadata freeze,
and exact affiliation wording remain deferred. Funding, conflicts of
interest, CRediT roles, data/code release terms, AI-use disclosure, and
venue-specific institutional checks also remain pending. No AI-generated
text in this shell constitutes author approval or submission
authorization.

\protect\phantomsection\label{refs}
\begin{CSLReferences}{1}{1}
\bibitem[\citeproctext]{ref-agarwal2025ttles}
Agarwal, Amit, Kushal Babel, Sourav Das, and Babak Poorebrahim
Gilkalaye. 2025. \emph{Time-Lock Encrypted Storage for Blockchains}.
Cryptology {ePrint} Archive, Paper 2025/2048.
\url{https://eprint.iacr.org/2025/2048}.

\bibitem[\citeproctext]{ref-eigenai2026}
Alves, David Ribeiro, Vishnu Patankar, Matheus Pereira, Jamie Stephens,
Nima Vaziri, and Sreeram Kannan. 2026. {``{EigenAI: Deterministic
Inference, Verifiable Results}.''} \emph{arXiv Preprint
arXiv:2602.00182}. \url{https://arxiv.org/abs/2602.00182}.

\bibitem[\citeproctext]{ref-bebel2022ferveo}
Bebel, Joseph, and Dev Ojha. 2022. \emph{Ferveo: Threshold Decryption
for Mempool Privacy in {BFT} Networks}. Cryptology {ePrint} Archive,
Paper 2022/898. \url{https://eprint.iacr.org/2022/898}.

\bibitem[\citeproctext]{ref-boneh2006threshold}
Boneh, Dan, Xavier Boyen, and Shai Halevi. 2006. {``Chosen Ciphertext
Secure Public Key Threshold Encryption Without Random Oracles.''}
\emph{Topics in Cryptology --- {CT-RSA} 2006}, 226--43.
\url{https://doi.org/10.1007/11605805_15}.

\bibitem[\citeproctext]{ref-boneh2025context}
Boneh, Dan, Benedikt Bünz, Kartik Nayak, Lior Rotem, and Victor Shoup.
2025. \emph{Context-Dependent Threshold Decryption and Its
Applications}. Cryptology {ePrint} Archive, Paper 2025/279; extended
abstract in ASIACRYPT 2025. \url{https://eprint.iacr.org/2025/279}.

\bibitem[\citeproctext]{ref-boneh2023accountability}
Boneh, Dan, Aditi Partap, and Lior Rotem. 2023. \emph{Accountability for
Misbehavior in Threshold Decryption via Threshold Traitor Tracing}.
Cryptology {ePrint} Archive, Paper 2023/1724.
\url{https://eprint.iacr.org/2023/1724}.

\bibitem[\citeproctext]{ref-bormet2025epochless}
Bormet, Jan, Sebastian Faust, Hussien Othman, and Ziyan Qu. 2025.
{``{BEAT-MEV}: Epochless Approach to Batched Threshold Encryption for
{MEV} Prevention.''} \emph{34th USENIX Security Symposium (USENIX
Security 25)}.
\url{https://www.usenix.org/conference/usenixsecurity25/presentation/bormet}.

\bibitem[\citeproctext]{ref-cerulli2023vetkeys}
Cerulli, Andrea, Aisling Connolly, Gregory Neven, Franz-Stefan Preiss,
and Victor Shoup. 2023. \emph{{vetKeys}: How a Blockchain Can Keep Many
Secrets}. Cryptology {ePrint} Archive, Paper 2023/616.
\url{https://eprint.iacr.org/2023/616}.

\bibitem[\citeproctext]{ref-chaum1992dleq}
Chaum, David, and Torben Pryds Pedersen. 1993. {``Wallet Databases with
Observers.''} \emph{Advances in Cryptology --- {CRYPTO} '92}, 89--105.
\url{https://doi.org/10.1007/3-540-48071-4_7}.

\bibitem[\citeproctext]{ref-choudhuri2025practical}
Choudhuri, Arka Rai, Sanjam Garg, Guru Vamsi Policharla, and Mingyuan
Wang. 2025. {``Practical Mempool Privacy via One-Time Setup Batched
Threshold Encryption.''} \emph{34th {USENIX} Security Symposium
({USENIX} Security 25)}, 3477--95.
\url{https://www.usenix.org/conference/usenixsecurity25/presentation/choudhuri}.

\bibitem[\citeproctext]{ref-rfc9380}
Faz-Hernandez, Armando, Sam Scott, Nick Sullivan, Riad S. Wahby, and
Christopher A. Wood. 2023. \emph{Hashing to Elliptic Curves}. RFC 9380.
Internet Engineering Task Force. \url{https://doi.org/10.17487/RFC9380}.

\bibitem[\citeproctext]{ref-herzberg1995proactive}
Herzberg, Amir, Stanislaw Jarecki, Hugo Krawczyk, and Moti Yung. 1995.
{``Proactive Secret Sharing or: How to Cope with Perpetual Leakage.''}
\emph{CRYPTO}, 339--52.

\bibitem[\citeproctext]{ref-lewis2025accountable}
Lewis-Pye, Andrew, Joachim Neu, Tim Roughgarden, and Luca Zanolini.
2025. {``Accountable Liveness.''} \emph{arXiv Preprint
arXiv:2504.12218}. \url{https://arxiv.org/abs/2504.12218}.

\bibitem[\citeproctext]{ref-liu2026uqvote}
Liu, Jingyu, Jingjing Rao, and Tetsutaro Uehara. 2026. {``{UQ-Vote}:
Consensus-Authorized Threshold Decryption for End-to-End Verifiable
Elections.''} \emph{IEEE Access} 14: 128905--16.
\url{https://doi.org/10.1109/ACCESS.2026.3725080}.

\bibitem[\citeproctext]{ref-eip7594}
Ryan, Danny, Dankrad Feist, Francesco D'Amato, Hsiao-Wei Wang, and Alex
Stokes. 2024. \emph{{EIP-7594: PeerDAS -- Peer Data Availability
Sampling}}. Ethereum Improvement Proposal.
\url{https://eips.ethereum.org/EIPS/eip-7594}.

\bibitem[\citeproctext]{ref-schoenmakers1999pvss}
Schoenmakers, Berry. 1999. {``A Simple Publicly Verifiable Secret
Sharing Scheme and Its Application to Electronic Voting.''}
\emph{CRYPTO}, 148--64. \url{https://doi.org/10.1007/3-540-48405-1_10}.

\bibitem[\citeproctext]{ref-shamir1979share}
Shamir, Adi. 1979. {``How to Share a Secret.''} \emph{Communications of
the ACM} 22 (11): 612--13.

\bibitem[\citeproctext]{ref-stengele2021ethtid}
Stengele, Oliver, Markus Raiber, Jörn Müller-Quade, and Hannes
Hartenstein. 2021. {``{ETHTID}: Deployable Threshold Information
Disclosure on {Ethereum}.''} \emph{2021 3rd International Conference on
Blockchain Computing and Applications ({BCCA})}, 127--34.
\url{https://doi.org/10.1109/BCCA53669.2021.9657019}.

\bibitem[\citeproctext]{ref-yurek2023dpss}
Yurek, Thomas, Zhuolun Xiang, Yu Xia, and Andrew Miller. 2023. {``Long
Live the Honey Badger: Robust Asynchronous {DPSS} and Its
Applications.''} \emph{32nd USENIX Security Symposium (USENIX Security
23)}, 5413--30.
\url{https://www.usenix.org/conference/usenixsecurity23/presentation/yurek}.

\end{CSLReferences}

\end{document}
